\section{Foundations and Related Work}

\subsection{Traffic digital twins for decision support}
\begin{bluenote}
Define the concept of a traffic digital twin in the context of urban traffic planning. It reviews
how digital twins are used for analyses and positions the work specifically on disruption management
and detour assessment.
\end{bluenote}

\subsubsection{Digital twins in traffic simulation}
\label{dt_traffic_simulation}

\textbf{Digital twin and data mining.} Recent advances have seen digital twin leveraged in urban
mobility to integrate real-time and historical data (from sensors, cameras and vehicles)
\cite{data_mining_dt}. The integration of advanced data mining techniques such as trajectory mining,
which exploits the information of trajectories for both private and public vehicles, and process
mining, which extracts processes from available event logs in information systems, within digital
twin frameworks can reveal mobility patterns, detect traffic anomalies, and simulate possible
scenarios.

However, the deployment of such systems is not without significant challenges. Foremost among these
is the need for continuous, reliable data acquisition. In addition, although mining trajectory data
enhances the ability to simulate mobility behaviors, it also raises significant privacy issues. For
large metropolitan areas, the demand for real-time processing introduces further complications, as
scenario simulations often need to meet stringent latency constraints.

\textbf{Digital twins in intelligent transport systems.} Digital twins in intelligent transport
systems (ITS) are digital replicas of real-world transport networks, providing a testbed for
simulating, analyzing, and optimizing traffic management strategies before their deployment in the
physical world \cite{intelligent_transport}. ITS integrate services across the city (including
traffic authorities, emergency responders, and public services) to enable greater coordination,
faster response times, and improved interactions between road users compared to conventional
systems. As discussed in \cite{intelligent_transport}, digital twins can model dynamics of urban
traffic flows by integrating real-time data from sensors, vehicles, and control centers, thereby
supporting predictive analytics and data-driven decision making.

These digital representations play a vital role in modern intelligent transportation systems,
enabling traffic management centers to shift from operator-dependent decision-making to more
automated and data-driven strategies. According to Rudskoy et al. \cite{intelligent_transport}, the
definition of a digital twin in this context is ``a module that reproduces a detailed digital model
of the road and allows for modeling and experiments.'' The primary purpose is to test solutions and
simulate various traffic scenarios before implementation in the real world.

The implementation of digital twins in ITS helps address key challenges in transport networks
including traffic optimization, management of transport during accidents or special events and
improving road safety.

\subsubsection{Simulation and mobility}

Simulation plays a pivotal role in mobility research, enabling the testing and analysis of various
traffic management and vehicle technologies in controlled, reproducible environments. Two notable
applications from recent literature are highlited below:

\textbf{AI-driven traffic simulation.} Talha Azfar and Ruimin Ke \cite{ai_driven_simulation} have
demonstrated the effectiveness of combining artificial intelligence with co-simulation environments
for urban mobility. Co-simulation means running different simulators together to model how connected
systems work as a whole \cite{cosimulation}. AI-agents using computer vision
(YOLO V5/V8\thesisurlfootnote{https://docs.ultralytics.com/fr/yolov5/}, which are models designed
for object detection) for real-time traffic observation were compared against statically configured
traffic lights. Performance metrics such as mean speed, total travel time and mean waiting time were
evaluated under medium and heavy traffic conditions. The results shows that the AI-agents reduce
mean waiting time and total trip duration especially with ground-truth data.

These results underscore the importance of simulation-based optimization. The integration of
multi-agent reinforcement learning and computer vision in a co-simulation framework enables
realistic evaluation and adaptation of traffic control strategies before real-world deployment.

\textbf{Simulation-based traffic light optimization.} Mehdi et al.
\cite{traffic_light_simulation} proposed an innovative modification to the traditional three-phase
traffic light system by introducing a four-stroke traffic light and directly measuring its impact at
a critical intersection in downtown Urmia, Iran. While previous research has explored a variety of
advanced algorithms (including queuing theory, stochastic modeling, variable-phase timing, and even
neural networks) to minimize vehicle waiting times, this work focuses on an alternative signal
configuration.

By using Arena Simulation Software\thesisurlfootnote{https://www.rockwellautomation.com/fr-be/products/software/arena-simulation.html},
the authors developed a comprehensive model capturing real-world vehicle arrival and departure
patterns at the intersection. Through systematic comparison of the conventional three-phase setup
with the proposed four-phase system across multiple traffic scenarios, the results indicated that
the four-stroke signal configuration significantly reduced the average time vehicles spent in the
system.

Although the study is centered on a single intersection, its findings highlight the potential
benefits of rethinking conventional traffic light algorithms. Future research should investigate
whether implementing such four-phase signal schemes at a broader, network-wide scale could further
alleviate congestion and improve overall urban mobility.

\subsubsection{Positioning of this work}

This work builds upon and extends the several existing contributions discussed before in the field
of traffic simulation and digital twins. A tutorial for traffic modeling using SUMO provides a
step-by-step guide to generate and simulate traffic scenarios \cite{sumo_tutorial}. This
methodology serves as the foundation for constructing a realistic simulation scenario for the
Brussels region, contributing to the development of a reliable digital twin based on local data.
Additionally, while a case study already exists for traffic simulation in Brussels using incomplete
data \cite{sparse_data_SUMO}, the proposed approach further leverages the HybridIoT methodology
introduced in that work to enhance model accuracy by estimating missing traffic information.
HybridIoT is a multi-agent system (MAS)-based approach designed to estimate missing environmental
data in large-scale environments \cite{hybridiot}. Lastly, as discussed before, an interactive tool
named \textit{TrafficTwin} has been developed to support traffic management experts in designing and
evaluating road deviation plans \cite{SIM_param}. The current work builds on this concept, adapting
and extending it. A new module is proposed for SUMO to facilitate the creation and simulation of
detour strategies, with the aim of assisting authorities such as
\emph{Bruxelles Mobilité}\thesisurlfootnote{https://be.brussels/fr/propos-de-la-region/structure-et-organisation/administrations-et-institutions-de-la-region/bruxelles-mobilite}
in optimizing roadwork planning. The overarching goal is to ensure the tool remains accessible, with
the potential for broader deployment to the general public.

\subsection{Traffic simulation for oprational planning}
\begin{bluenote}
Introduce microscopic traffic simulation principles and explains why SUMO is appropriate for the
work. It also discusses key limitations.
\end{bluenote}

\subsubsection{Traffic simulation tools}

Traffic congestion is one of the major challenges faced by modern cities and is expected to get
worse with increasing population and vehicle usage \cite{increasing_traffic,increasing_cars}. Traffic
simulation tools have been developed as mathematical and computational models that replicate
real-world traffic behavior and dynamics \cite{TST_def}. These tools help transportation planners
and engineers to optimize road networks designs or estimate traffic-related emissions.

Traffic simulations are commonly categorized into three main types based on the level of detail
desired \cite{TST_types}:
\begin{enumerate}
\item \textbf{Microscopic Simulation}: Models each traffic participant (e.g., vehicles,
pedestrians) individually. Agents continuously in time and space respond to their immediate
environment. The global traffic dynamics is the result of the sum of these individual decisions and
behaviors.
\item \textbf{Mesoscopic Simulation}: Still represents individual agents, but their behavior is
determined by average traffic conditions like speed or density. Agents only interact directly at key
points in the network, such as intersections.
\item \textbf{Macroscopic Simulation}: Looks at traffic as a whole rather than focusing on
individual vehicles.
\end{enumerate}

Choosing a simulation type entails balancing computational cost and realism. For instance,
microscopic simulation is ideal for projects that require modeling of individual behaviors or the
testing of interventions at the street or intersection level, which is often necessary for a dense
urban environment such as Brussels. In contrast, a macroscopic simulation might suffice for
citywide traffic forecasting or strategic planning, where the focus is on broader patterns rather
than local details.

\subsubsection{SUMO}

\textbf{SUMO} (Simulation of Urban MObility) is an open-source, highly portable, microscopic and
continuous multimodal traffic simulation tool. It allows detailed modeling of individual vehicles,
each following its own route through a specified road network. As a microscopic simulator, SUMO
treats every vehicle as an individual entity with specific behaviors and interactions. The tool is
multi-modal enabling the simulation of diverse traffic participants including cars, public transport
vehicles and pedestrians \cite{SUMO_eclipse,SUMO_homepage,SUMO_doc_glance}.

SUMO provides some tools to setup the simulation and managing the data of the road network
\cite{SUMO_eclipse,SUMO_doc_glance}. In this project, the following components of the SUMO
ecosystem are particularly relevant:
\begin{itemize}
\item \textbf{sumo}: Executes command-line-based microscopic traffic simulations without graphical
visualization.
\item \textbf{netconvert}: Imports and converts road network data from various formats into the
internal SUMO network format.
\item \textbf{duarouter}: Computes optimal routes across the network based on demand inputs and
performs Dynamic User Assignment (DUA).
\end{itemize}

SUMO has been used extensively to model real-world traffic systems in diverse urban contexts,
including cities such as Berlin, Dublin and Turin \cite{SUMO_scenarios}. Each scenario adapts to
the specific traffic dynamics and urban characteristics of its location:

\begin{enumerate}
\item \textbf{Berlin}: Berlin has served for large-scale traffic simulations using SUMO. The ``24
Hours of Traffic in Berlin'' scenario captures over 2.2 million trips within an 800km² area. This
model enhances traditional traffic simulations by incorporating motorized private vehicles alongside
realistic bus schedules derived from actual General Transit Feed Specification (GTFS) data,
representing 89\% of the city's bus lines. This scenario allows for holistic analyses that combine
traffic, communication (including V2X and LTE/5G) and application simulations, making it suited for
smart mobility researches. \cite{berlin_scenario}
\item \textbf{Dublin}: The study by Gueriau and Dusparic \cite{dublin_scenario} leverages SUMO to
simulate and analyze the impact of connected and autonomous vehicles (CAVs) on real Irish road
networks, covering urban, national and motorway settings. Using detailed demand and network data
from Dublin and surrounding areas, the research quantifies how varying CAV penetration rates affect
both traffic efficiency and open dataset for benchmarking the effects of emerging vehicle
technologies in mixed traffic environments.
\item \textbf{Turin}: Leveraging real-world data, the TuSTScenario \cite{turin_scenario} provides a
highly detailed SUMO-based simulation of metropolitan Turin, covering over 602.61km². The model
integrates structural road information from OpenStreetMap, validated average speeds from Google
APIs, actual traffic light programming and real trip demands from local mobility services. With more
than 2.2 million simulated daily trips and validated against city sensor data, this scenario enables
realistic studies of urban traffic dynamics and congestion patterns.
\end{enumerate}

\textbf{Brussels} is also a large city warranting similarly large-scale simulations. However, the
focus of this study is distinct: rather than prioritizing multi-modal integration or CAV impact, we
specifically investigate the planning and management of roadworks within the Brussels network. The
goal is to demonstrate how such interventions can be realistically modeled, analyzed, and optimized
in SUMO for modern urban infrastructure management.

Beyond traditional traffic flow analysis SUMO has been used to evaluate the effects of technologies
like virtual traffic lights communication systems where cars directly communicate with each other
\cite{SUMO_VTL}. This study showed that replacing traffic lights with virtual ones can reduce both
congestion and the likelihood of traffic accidents. Moreover SUMO can be used in environmental
research, for instance in assessing the impact of traffic-related emissions on urban air quality
\cite{SUMO_emission}.

One of the strengths of SUMO is its deterministic core which allows reproducibility of experiments.
However, this deterministic nature can be a limitation when it comes to capturing the variability in
human driving behavior which requires the integration of stochastic elements \cite{SUMO_sumo}.

\subsubsection{Alternatives to SUMO}

Several other traffic simulation tools can serve as alternatives to SUMO, each with different
features, advantages and limitations:
\begin{enumerate}
\item \textbf{PTV Vissim}: PTV-Vissim\thesisurlfootnote{https://www.ptvgroup.com/en/products/ptv-vissim}
is a widely used microscopic traffic simulation software that also supports multi-modal transport
systems. It offers a broad set of features comparable to those of SUMO. However, it is a
proprietary tool that requires a commercial license, which can restrict its accessibility and limit
the ability to implement custom modifications at the code level.
\item \textbf{Aimsun}: Aimsun\thesisurlfootnote{https://docs.aimsun.com/} is a commercial traffic
simulation platform that supports microscopic, mesoscopic, and macroscopic modeling. While it is
used in various contexts, its commercial nature introduces licensing costs and reduces flexibility
for developers who require deeper code-level access and customization.
\item \textbf{CityFlow}: CityFlow \cite{cityflow} an open-source microscopic simulator designed for
efficient, large-scale urban traffic simulations, capable of handling thousands of intersections
with high computational efficiency. Its architecture is highly optimized for speed and parallel
execution, making it over twenty times faster than SUMO in large-scale benchmarks
\cite{cityflow_rl}. CityFlow is particularly popular in the field of artificial intelligence and
reinforcement learning, serving as a benchmark for multi-agent adaptive traffic signal control and
dynamic routing tasks. However, CityFlow remains limited compared to SUMO in certain areas: it
offers fewer out-of-the-box features for conventional traffic simulation, has a less mature
ecosystem, and does not natively support the simulation of public transportation modes such as
buses, trams, or trains, which restricts its suitability for comprehensive, multi-modal urban
simulations.
\end{enumerate}

\subsubsection{Why prefer SUMO}

SUMO was selected for this project based on the following criteria:
\begin{itemize}
\item \textbf{Open-source}: SUMO's open-source status allows users to freely modify, adapt, and
extend the software, making it suitable for research and academic use where flexibility and
transparency are important. This degree of code-level access is not available in proprietary tools
like PTV Vissim or Aimsun.
\item \textbf{Community and documentation}: SUMO benefits from a large active community that
contributes regularly to new features or bug fixes. In addition, it offers extensive documentation
and a wide range of tutorials, which significantly lower the learning curve for new users.
\item \textbf{Microscopic and multi-modal simulation}: SUMO provides detailed agent-based
(microscopic) modeling of cars, public transport, and pedestrians. While other simulators such as
Vissim and Aimsun also support microscopic and multi-modal simulation, SUMO's implementation is
open, which enables deeper customization and integration for research needs.
\item \textbf{Integration and flexibility}: SUMO's design and command-line tools (like
\texttt{netconvert} and \texttt{duarouter}) facilitates working on various data sources making it
easy to connect with external components.
\end{itemize}

Commercial tools like PTV Vissim and Aimsun offer advanced graphical user interfaces and support;
however, their proprietary nature limits customizability and long-term usability in research
settings. CityFlow, while open-source and scalable, lacks SUMO's extensive documentation and wider
range of traffic simulation features relevant for this project's objectives.

\subsection{Metrics for urban traffic planning}
\begin{bluenote}
Explain what ``impact'' is, it can be defined at different levels, network efficiency, specific user
impact, delay etc. Motivates the need of a multi-metric perpective and sets the foundations for
later scoring and ranking strategies
\end{bluenote}

\subsection{Stochastic simulation}
\begin{bluenote}
Why a single run is not enough, introduce Monte Carlo simulation, estimating performance
distributions and supporting decisions with confidence.
\end{bluenote}

In the context of a Digital Twin for traffic management (especially one designed to evaluate
deviation plans for disruptions like roadworks or accidents) the treatment of uncertainty is not
merely a technical refinement but a fundamental requirement for validity
\cite{digital_twin_transport,flow_twin}. Traffic is a stochastic process and small perturbations in
the initial conditions can propagate to cause macroscopic differences in flow and congestion
patterns \cite{stochastic_macroscopic}.

Stochastic simulation introduces randomness into the model to mimic the real-world uncertainties in
demand, driver behaviour and network conditions \cite{application_stochastic_risk}. By treating
model inputs and parameters as random variables rather than fixed values, stochastic simulations
generate a distribution of possible outcomes that better reflects the variability of traffic states.
This transition from point-estimates to probabilistic distributions allows for the application of
risk assessment frameworks, enabling planners to weigh the trade-offs between efficiency and
reliability.

\subsubsection{Theoritical foundations of the Monte Carlo method}

The Monte Carlo method is a broad class of computational algorithms that rely on repeated random
sampling to obtain numerical results. The underlying concept is to use randomness to solve problems
that might be deterministic in principle but are too complex to solve analytically due to the high
dimensionality of the problem.

\subsection{Multiple-criteria decision analysis}

The multiple-criteria decision analysis (MCDA) aims to help in situations where multiple criteria,
most of the time contradictory, must be.
